\newif\ifglobal

%%%%%%%%%%%%%%%%%%%%%%%
%%% Compile to view %%%
%%%%%%%%%%%%%%%%%%%%%%%

%%%%%%%%%%%%%%%%%%%%%%%%%%%%%%%%%%%%%%%%%%%%%%%%%%%
%%% Readme:                                     %%%
%%% https://www.overleaf.com/read/bwdjvfjksvjy/ %%%
%%%%%%%%%%%%%%%%%%%%%%%%%%%%%%%%%%%%%%%%%%%%%%%%%%%

% >>> M?: marks a module setting number or important notice

% >>> M4: Set {./relative-path} to external files for this module <<<
\def \modulefiles {./38-TWAI}
\newcommand{\modulename}{CAN FD}
% To add an external file, use:
% (Replace '---' with actual filename)
% '\input{\modulefiles/tables/---.tex}' for tables
% '\includegraphics{\modulefiles/figures/---}' for figures
% '\subfile{\modulefiles/---}' for register files

% >>> M1: This module is in global project? <<<
% 'YES' - uncomment; 'NO' - comment out
\globaltrue

\ifglobal
    \documentclass[main\_\_CN.tex]{subfiles}

\else
    % Font "FandolSong-Regular" used by ctex does not have GSUB table.
    % It causes warnings that do not affect compilation and can be ignored.
    \PassOptionsToPackage{quiet}{fontspec}

    \documentclass[openany, 10pt]{book}

    % >>> M2: In ./00-shared/preamble-trm-module, also find
    %         settings for visibility of todo notes, watermark <<<
    \usepackage{preamble-trm-module}

    % Enable Chinese typesetting
    \usepackage{ctex}

    % >>> M3: Temporary labels in Chinese
    %         you can add your own labels there <<<
    \myexternaldocument{./00-shared/temp-labels-cn}

\fi %\ifglobal

\setcounter{secnumdepth}{4}

% Define formatting of chapters in text
\titleformat{\chapter}
  [display]
  {\LARGE\bfseries}
  {\color{AllportsColor}第\, \thechapter\, 章}
  {1em}
  {\color{AllportsColor}}
\titlespacing*{\chapter}{0pt}{0pt}{20pt}

% Define headers
\usepackage{fancyhdr}
\usepackage{lipsum}
\renewcommand{\chaptermark}[1]{\markboth{第\, \thechapter\, 章 \  \ #1}{}}
\fancypagestyle{plain}{%
  \fancyhead{}%
  %\fancyhead[R]{\Navigation}
  \fancyhead[L]{%
    \nouppercase{%
      \ %
      \normalsize\leftmark
    }%
  }%
}
\pagestyle{plain}

% >>> M5: If this module needs any updates in
% './00-shared/preamble-trm-module.sty'
% (include additional package, etc.), contact your module PM
% or create the file './00-shared/changelog.tex' make the
% necessary updates yourself and document them in this file <<<

% Make text left-aligned and allow latex to break pages naturally, comment out for full justification
\raggedright
\raggedbottom

\begin{document}

% Sets language into which pre-defined bits of text are translated
\selectlanguage{chinese}

% Do not delete this block
\ifglobal
% Add the following front matter if
% this module is outside of global project
\else
    \InputIfFileExists{readme}{}{}
    \listoftodos
    {\let\clearpage\relax \listofdonetodos}
    \tableofcontents
    %\listoftables
    %\listoffigures
\fi

%%%%%%%%%%%%%%%%%%%%%%%%%
%%% TRM Module Begins %%%
%%%%%%%%%%%%%%%%%%%%%%%%%

\hypertarget{canfd}{}
\chapter{CAN FD}
\label{mod:twai}


\modulename{}（Controller Area Network Flexible Data-Rate，控制器局域网灵活数据速率）协议是一种多主机、多播的通信协议，具有错误检测、发送错误信号以及内置报文优先仲裁等功能，遵循 ISO11898-1 规定的 CAN 协议。CAN 协议适用于汽车和工业应用。

\chipname{} 中集成的 CAN FD 控制器与 CAN FD 规范协议兼容，但尚未获得正式认证。除非另有说明，本章中提到的 CAN FD 控制器均指兼容 CAN FD 协议的控制器。


%%%%%%%%%%%%%%%%
%%% Overview %%%
%%%%%%%%%%%%%%%%

\section{概述}


\chipname{} 包含两个 \modulename{} 控制器，可以通过外部收发器连接到 CAN FD 总线。\modulename{} 支持许多高级功能，应用场景广泛，如汽车产品、工业自动化控制、楼宇自动化等。

本章节提供了 \modulename{} 的功能描述、程序模型和参数，旨在为软件驱动开发人员提供参考。


%%%%%%%%%%%%%%%%%%%%
%%% Feature List %%%
%%%%%%%%%%%%%%%%%%%%

\section{主要特性}
\label{sec:twai-features}

\chipname{} 的 \modulename{} 具有以下特性：

\begin{itemize}
    \item 兼容 ISO11898-1:2015 标准
    \item 具有 128 字的 RX FIFO（6 个 64 字节载荷的 CAN FD 帧，21 个 8 字节载荷的 CAN/CAN FD 帧）
    \item 4 个 TX buffer（每个 TX buffer 中有 1 个 CAN FD 帧）
    \item 32 位 APB 配置接口
    \item 支持 ISO 和非 ISO CAN FD 协议
    \item 支持时间戳和定时发送
    \item 中断
    \item 内置 3 个屏蔽过滤器和 1 个范围过滤器
    \item 回环模式、总线监控模式、禁止确认模式、自检模式、受限操作模式
\end{itemize}


%%%%%%%%%%%%%%%%%%%%%%%%%%%%%%
%%% Functional Description %%%
%%%%%%%%%%%%%%%%%%%%%%%%%%%%%%

\section{功能描述}
\label{sec:twai-func-descr}

\subsection{时钟}
\label{sec:twai-func-clk}

\modulename{} 采用单一时钟设计，唯一的时钟是 TWAI\_CLK，所有时间参数均由此派生。

\begin{figure}[H]
    \includegraphics[width=0.9\linewidth]{\modulefiles/figures/esp32c5-twai-clk.png}
    \centering
    \caption{TWAI\_CLK Diagram}
    \label{fig:twai-clk}
\end{figure}

\chipname{} \modulename{} 支持 XTAL\_CLK 与 PLL\_F80M\_CLK 两种时钟源选择：

\begin{itemize}
    \item XTAL\_CLK (外部晶振)：高精度时钟源，支持最高频率为 48 MHz。
    \item PLL\_F80M\_CLK (锁相环时钟)：倍频时钟源，最高工作频率可达 80 MHz。
\end{itemize}

可通过 \hyperref[fielddesc:PCRTWAI0FUNCCLKSEL]{PCR\_TWAI0\_FUNC\_CLK\_SEL} 和 \hyperref[fielddesc:PCRTWAI1FUNCCLKSEL]{PCR\_TWAI1\_FUNC\_CLK\_SEL} 寄存器选择两个 CAN FD 控制器的时钟源。

\begin{itemize}
    \item 当数据速率 $\le$ 5 Mbit/s 时：推荐选用 XTAL\_CLK，以获得更高的频率稳定度。
    \item 当数据速率 > 5 Mbit/s 时：推荐选用 PLL\_F80M\_CLK，以满足高带宽需求。
\end{itemize}

具体的位速率（相关介绍详见 \ref{sec:twai-func-bit-rate} \textit{\nameref{sec:twai-func-bit-rate}}）设定，请参考 \hyperref[regdesc:TWAIFDBTRFDREG]{TWAIFD\_BTR\_FD\_REG} 和 \hyperref[regdesc:TWAIFDBTRREG]{TWAIFD\_BTR\_REG} 寄存器的参数计算说明。

\subsection{复位}


上电后，\modulename{} 应通过硬件或软件进行复位。将 1 写入 \hyperref[fielddesc:TWAIFDRST]{TWAIFD\_RST} 可以实现软件复位。如需硬件复位，则 \modulename{} 在系统时钟的两个时钟周期内不可访问。例如，如果 \modulename{} 系统时钟为 100 MHz，软件应在硬件复位释放后等待 20 ns 再访问 \modulename{}。如需软件复位，则不必等待。硬件复位和软件复位作用相同。无论采用哪种复位方式，\modulename{} 都将进入以下状态：

\begin{itemize}
    \item 禁用 \modulename{}，不能在 CAN 总线上通信（离线状态）。
    \item \modulename{} 所有内存寄存器取复位值。
    \item \modulename{} 内存（TX buffer 和 RX buffer）未复位。
\end{itemize}




\subsection{时间基准}
\label{sec:twai-func-time-base}



\modulename{} 集成了独立的 32 位定时器，可以产生时间戳，为定时发送或接收帧的提供时间基准。时间基准定时器可以通过配置 \hyperref[fielddesc:TWAIFDTIMERCE]{TWAIFD\_TIMER\_CE} 启用，然后通过 \hyperref[fielddesc:TWAIFDTIMESTAMPHIGH]{TWAIFD\_TIMESTAMP\_HIGH} 和 \hyperref[fielddesc:TWAIFDTIMESTAMPLOW]{TWAIFD\_TIMESTAMP\_LOW} 读取其当前值。

时间基准是一个向上计数的无符号定时器，\hyperref[fielddesc:TWAIFDTIMERSTEP]{TWAIFD\_TIMER\_STEP} 配置计时步长，\hyperref[fielddesc:TWAIFDTIMERLDVALL]{TWAIFD\_TIMER\_LD\_VAL\_L} 配置定时器的初始值，\hyperref[regdesc:TWAIFDTIMERCTVALLREG]{TWAIFD\_TIMER\_CT\_VAL\_L} 配置定时器的终到值。时间基准的有效位数可通过 \hyperref[fielddesc:TWAIFDTSBITS]{TWAIFD\_TS\_BITS} 寄存器读取或配置，有效范围为 1 到 32。

\subsection{工作模式}
\label{sec:twai-func-oper-mode}

复位后，\modulename{} 处于禁用状态，不参与 CAN 总线通信（不发送、接收或监控数据）。在启用 \modulename{} 之前，必须按照章节 \ref{sec:twai-func-bus-config} \textit{\nameref{sec:twai-func-bus-config}} 中的说明完成必要配置。配置完成之后，对 \hyperref[fielddesc:TWAIFDENA]{TWAIFD\_ENA} 写 1 启用 \modulename{}。当 \hyperref[fielddesc:TWAIFDENA]{TWAIFD\_ENA} = 1 时，\modulename{} 开始接入 CAN 总线，并在接收到 11 个连续的隐性位后正式加入 CAN 总线通信。加入总线后，\modulename{} 进入主动错误状态（在接入过程中处于离线状态），开始在 CAN 总线上进行数据交互。\modulename{} 加入 CAN 总线通信的时刻可以通过 FCS 中断以及后续读取 \hyperref[regdesc:TWAIFDEWLERPFAULTSTATEREG]{TWAIFD\_EWL\_ERP\_FAULT\_STATE\_REG} 来确定（详见章节 \ref{sec:twai-interrupt} \textit{\nameref{sec:twai-interrupt}}）。\modulename{} 的基本工作模式如图 \ref{fig:twai-oper-mode} 所示。

\begin{figure}[H]
\begin{centering}
\includegraphics{\modulefiles/figures/FSM_main_modes}
\par\end{centering}
\caption{工作模式}
\label{fig:twai-oper-mode}
\end{figure}

\modulename{} 处于主动错误状态时，仍然参与 CAN 总线通信。如果 \modulename{} 变为被动错误状态并随后进入离线状态，则会停止在 CAN 总线上通信，直至收到错误计数器复位命令（向 \hyperref[fielddesc:TWAIFDERCRST]{TWAIFD\_ERCRST} 写 1）。收到此命令后，\modulename{} 开始重新接入总线，直到 \modulename{} 检测到 129 组由 11 个连续隐性位组成的序列（实现层面的选择，超过 CAN 规范规定的最小值 128 个序列），之后再次进入主动错误状态。

清零 \hyperref[fielddesc:TWAIFDENA]{TWAIFD\_ENA} 寄存器，可以随时禁用 \modulename{}。在这种情况下：

\begin{itemize}
    \item \modulename{} 立即停止在 CAN 总线上通信，并仅发送隐性位。
    \item TX/RX 错误计数器 (TEC/REC) 复位为 0，\modulename{} 进入离线状态。
    \item 所有 TX buffer 切换至 empty 状态（见图 \ref{fig:twai-tx-buffer-state}），TX buffer RAM 中的内容仍然有效（内存未复位）。
    \item RX buffer 刷新（详见章节 \ref{sec:twai-func-flush} \textit{\nameref{sec:twai-func-flush}}）。
\end{itemize}

建议在清零 \hyperref[fielddesc:TWAIFDENA]{TWAIFD\_ENA} 禁用 \modulename{} 之前，避免发送任何帧，防止 CAN 总线上其他节点发送错误帧。因此，操作 \modulename{} 的软件驱动程序应确保 \modulename{} 内的任何 TX buffer 不处于 ready、TX in progress 或 abort in progress 状态（详见章节 \ref{sec:twai-can-frame-trans} \textit{\nameref{sec:twai-can-frame-trans}}）。

\hyperref[fielddesc:TWAIFDERCRST]{TWAIFD\_ERCRST} 是一个“粘性”命令。这意味着如果 \modulename{} 尚未关闭总线且此命令已发出，\modulename{} 将记录此命令，在进入离线状态的瞬间自动触发重新接入总线。这样，软件驱动程序可以在通信期间提前发送该命令，从而在离线后立即触发重新接入，避免额外的软件延迟。


\subsection{初始化顺序}

初始化顺序包括以下步骤：

\begin{enumerate}
    \item 复位（硬件或软件复位）
    \item 配置 \modulename{}：
    \begin{enumerate}
        \item 按照章节 \ref{sec:twai-interrupt} \textit{\nameref{sec:twai-interrupt}} 配置中断
        \item 按照章节 \ref{sec:twai-func-bit-rate} \textit{\nameref{sec:twai-func-bit-rate}} 配置位速率
        \item 配置过滤器、特殊模式等其他功能
    \end{enumerate}
    \item 对 \hyperref[fielddesc:TWAIFDENA]{TWAIFD\_ENA} 写 1 启用 \modulename{}。
    \item 轮询 \hyperref[regdesc:TWAIFDEWLERPFAULTSTATEREG]{TWAIFD\_EWL\_ERP\_FAULT\_STATE\_REG} 寄存器，或等待故障限制状态触发中断 \hyperref[fielddesc:TWAIFDFCSIINTST]{TWAIFD\_FCSIINTST}。\hyperref[fielddesc:TWAIFDERA]{TWAIFD\_ERA} = 1（\modulename{} 进入主动错误状态）表示接入完成。
    \item 初始化完成，软件驱动程序可以发送和接收帧。
\end{enumerate}

\subsection{反初始化顺序}

反初始化顺序包括以下步骤：

\begin{enumerate}
    \item 确保没有任何 TX buffer 处于 ready、TX in progress 或 abort in progress 状态：向 TX buffer 发出 set abort 命令（详见章节 \ref{sec:twai-can-frame-trans} \textit{\nameref{sec:twai-can-frame-trans}}），并且避免将新的帧插入 TX buffer 进行传输。
    \item 清零 \hyperref[fielddesc:TWAIFDENA]{TWAIFD\_ENA}。
\end{enumerate}

\subsection{CAN 总线配置\label{sec:twai-func-bus-config}}

\subsubsection{位速率}
\label{sec:twai-func-bit-rate}

CAN 总线上的位速率来源于系统时钟（参见章节 \ref{sec:twai-func-clk} \textit{\nameref{sec:twai-func-clk}}）。CAN 总线上的基本时间单位是时间定额 (time quanta)，计算方式为将系统时钟频率除以位速率预分频器的值。\modulename{} 具有独立的预分频器用于标称位速率（通过 \hyperref[fielddesc:TWAIFDBRP]{TWAIFD\_BRP} 配置）和数据位速率（通过 \hyperref[fielddesc:TWAIFDBRPFD]{TWAIFD\_BRP\_FD} 配置）。配置位速率，可通过 \hyperref[regdesc:TWAIFDBTRREG]{TWAIFD\_BTR\_REG}（标称位速率）和 \hyperref[regdesc:TWAIFDBTRFDREG]{TWAIFD\_BTR\_FD\_REG}（数据位速率）寄存器指定 Prop\_Seg、Phase\_Seg1 和 Phase\_Seg2 的持续时间，如图 \ref{fig:twai-bit-rate} 所示。

\begin{figure}[H]
\begin{centering}
\includegraphics[scale=0.75]{\modulefiles/figures/Bit_time}
\par\end{centering}
\caption{位时间}
\label{fig:twai-bit-rate}
\end{figure}

配置流程详见章节 \ref{sec:twai-prog-500-4-example} \textit{\nameref{sec:twai-prog-500-4-example}}。


\subsubsection{发送器延迟}

发送器延迟是 \modulename{} 在 CAN\_TX 输出上发送的信号返回 CAN\_RX 输入的传播延迟，如图 \ref{fig:twai-trans-delay} 所示。此延迟包括信号传播到物理层收发器的延迟、收发器本身的延迟，以及从收发器到 CAN\_RX 输入的延迟。无论帧中是否切换位速率，\modulename{} 在发送 CAN FD 帧时会测量其自身的发送器延迟。测量位置在 FDF (EDL) 位和 R0 位之间由隐性到显性的边沿，如图 \ref{fig:twai-trans-delay-measure} 所示。发送器延迟可以通过 \hyperref[fielddesc:TWAIFDTRVDELAYVALUE]{TWAIFD\_TRV\_DELAY\_VALUE} 进行测量，其单位为系统时钟周期。

\begin{figure}[H]
    \centering
    \includegraphics[scale=0.9]{\modulefiles/figures/Transmitter_delay}
    \caption{发送器延迟}
    \label{fig:twai-trans-delay}
\end{figure}

\begin{figure}[H]
    \centering
    \includegraphics[scale=0.9]{\modulefiles/figures/Transmitter_delay_measurement}
    \caption{发送器延迟测量}
    \label{fig:twai-trans-delay-measure}
\end{figure}

测量的发送器延迟包括 \modulename{} 的输入延迟，即系统时钟的 2 个周期。因此，测量的发送器延迟将始终比从 CAN\_TX 到 CAN\_RX 的实际延迟多两个周期。例如，如果从 CAN\_TX 到 CAN\_RX 的信号传播需要 110 ns（在 100 MHz 时为 11 个系统时钟周期），测量的发送器延迟将为 13 个周期。

发送器延迟的测量有最大值限制，127 个系统时钟周期饱和。如果 CAN\_TX 和 CAN\_RX 之间的延迟超过此值，测量结果将为 127 个周期。当系统时钟频率为 100 MHz 时，这意味着最大可测量的发送器延迟为 1.27 $\mu s$，此值大于大多数 CAN 收发器所需的延迟。

\subsubsection{二次采样点}

二次采样点用于在数据位速率期间由 \modulename{} 检测比特错误。二次采样点的位置从位时间的开始 (Sync\_Seg) 起，表示为系统时钟周期的延迟（而不是时间定额）。二次采样点位置可以是固定的（通过 \hyperref[fielddesc:TWAIFDSSPOFFSET]{TWAIFD\_SSP\_OFFSET} 配置），从发送器延迟派生（通过 \hyperref[fielddesc:TWAIFDSSPOFFSET]{TWAIFD\_SSP\_OFFSET} 和 \hyperref[fielddesc:TWAIFDTRVDELAYVALUE]{TWAIFD\_TRV\_DELAY\_VALUE} 配置），或者可以禁用（无 SSP），如图 \ref{fig:twai-second-sampling-point} 所示。当禁用二次采样点时，\modulename{} 在数据位速率下发送时使用 hyperref[regdesc:TWAIFDBTRREG]{TWAIFD\_BTR\_REG} 配置的常规采样点。

\begin{figure}[H]
\begin{centering}
\includegraphics[scale=0.9]{\modulefiles/figures/SSP_logic}
\par\end{centering}
\caption{二次采样点}
\label{fig:twai-second-sampling-point}
\end{figure}

\begin{figure}[H]
\begin{centering}
\includegraphics[scale=0.9]{\modulefiles/figures/SSP_timing}
\par\end{centering}
\caption{二次采样点 2}
\label{fig:twai-second-sampling-point-2}
\end{figure}

二次采样点偏移量 (\hyperref[fielddesc:TWAIFDSSPOFFSET]{TWAIFD\_SSP\_OFFSET}) 可配置范围为 0～255，二次采样点位置的内部范围同为 0～255。如果出于某些原因计算出的二次采样点位置超过 255 个时钟周期（例如测量的发送器延迟较大），则其值最大为 255。

由于 \modulename{} 具有 2 个系统时钟周期的输入延迟（最小时间定额），建议二次采样点偏移量至少设置为 2 以补偿该延迟。如果 \hyperref[fielddesc:TWAIFDSSPOFFSET]{TWAIFD\_SSP\_OFFSET} < 3 且 \hyperref[fielddesc:TWAIFDSSPSRC]{TWAIFD\_SSP\_SRC} = SSP\_SRC\_OFFSET，则 \modulename{} 无法检测自身在 \modulename{} 发送帧的比特错误。

\modulename{} 在使用二次采样点时最多可以处理 CAN\_TX 和 CAN\_RX 管脚之间的 4 个飞行比特。例如，在系统时钟频率为 100 MHz、数据位速率为 5 Mbit/s 的情况下，一个数据位时间对应 20 个系统时钟周期。因此，二次采样点的最晚可能位置应设定为 4 × 20 = 80 个系统时钟周期。用户在配置二次采样点时，应确保其最终位置（包括 \hyperref[fielddesc:TWAIFDSSPOFFSET]{TWAIFD\_SSP\_OFFSET}/\hyperref[fielddesc:TWAIFDTRVDELAYVALUE]{TWAIFD\_TRV\_DELAY\_VALUE}）不超过 4 个数据位时间。


\subsubsection{CAN FD 支持}

\modulename{} 支持选择 ISO 和非 ISO 版本的 CAN FD 协议，可以通过 \hyperref[fielddesc:TWAIFDNISOFD]{TWAIFD\_NISOFD} 配置。选择 ISO 协议版本时，\modulename{} 符合 ISO11898-1:2015 规范。选择非 ISO 版本时，\modulename{} 符合 CAN FD 1.0 规范。只有在 \modulename{} 禁用时 (\hyperref[fielddesc:TWAIFDENA]{TWAIFD\_ENA} = 0)，才能配置 \hyperref[fielddesc:TWAIFDNISOFD]{TWAIFD\_NISOFD}。

\subsubsection{协议异常处理}

\modulename{} 支持协议异常检测，可以通过 \hyperref[fielddesc:TWAIFDPEX]{TWAIFD\_PEX} 启用。只有在 \modulename{} 禁用时 (\hyperref[fielddesc:TWAIFDENA]{TWAIFD\_ENA} = 0)，才能配置 \hyperref[fielddesc:TWAIFDPEX]{TWAIFD\_PEX}。

协议异常的具体行为取决于 CAN 的实现类型，详见表 \ref{tab:twai-func-can-implt-type}。当 \hyperref[fielddesc:TWAIFDPEX]{TWAIFD\_PEX} = 1 且 \modulename{} 检测到协议异常，模块将进入总线接入状态，并等待在 CAN\_RX 信号上连续检测到 11 个隐性位。

此过程中，协议异常不会影响 REC/TEC 计数器的值，\modulename{} 的故障限制状态也保持不变。检测到协议异常将会置位 \hyperref[fielddesc:TWAIFDPEXS]{TWAIFD\_PEXS} 标志，可通过对 \hyperref[fielddesc:TWAIFDCPEXS]{TWAIFD\_CPEXS} 写 1 来清除该标志。

若 \hyperref[fielddesc:TWAIFDPEX]{TWAIFD\_PEX} = 0 （即未启用协议异常检测）且协议异常仍然存在，\modulename{} 将发送错误帧。

\subsubsection{实现类型}

ISO11898-1:2015 定义了三种 CAN 协议实现类型：经典 CAN、CAN FD 容忍 (CAN FD tolerant) 和 CAN FD 启用 (CAN FD enabled)。\modulename{} 支持三种实现类型。通过 \hyperref[fielddesc:TWAIFDFDE]{TWAIFD\_FDE} 和 \hyperref[fielddesc:TWAIFDPEX]{TWAIFD\_PEX} 可以配置符合实现类型的协议，上述字段仅在 \modulename{} 禁用时 (\hyperref[fielddesc:TWAIFDENA]{TWAIFD\_ENA} = 0) 才能配置。

\begin{table}[H]
    \centering
    \caption{CAN 实现类型}
    \label{tab:twai-func-can-implt-type}
    \begin{tabular}{ |m{2.5cm}|m{1.5cm}|m{1.5cm}|m{9.2cm}| }
    \hline
    \rowcolor{lightgray}
    \textbf{实现类型} & \textbf{\hyperref[fielddesc:TWAIFDFDE]{TWAIFD\newline \_FDE}} & \textbf{\hyperref[fielddesc:TWAIFDPEX]{TWAIFD\newline \_PEX}} &   \textbf{行为} \\\hline

    经典 CAN & 0 & 0 & \modulename{} 检测到隐性 FDF 位（基本帧中 IDE 之后的位和扩展帧中 RTR/R1 之后的位）时，会发送错误帧。\\\hline
    CAN FD 容忍 & 0 & 1 & \modulename{} 检测到隐性 FDF 位时，将视为协议异常，并进入总线接入状态。\\\hline
    CAN FD 启用 & 1 & 0 & \modulename{} 能够接收和发送 CAN FD 帧。\modulename{} 在 res 位（FDF 之后的位）检测到隐性值时，会发送错误帧。\\\hline
    CAN FD 启用且协议异常 & 1 & 1 & \modulename{} 能够接收和发送 CAN FD 帧。\modulename{} 在 res 位（FDF 之后的位）检测到隐性值时，将视为协议异常，并进入总线接入状态。此配置兼容 CAN FD 协议的未来扩展（例如 CAN XL）。\\\hline
    \end{tabular}
\end{table}

当 \modulename{} 配置为经典 CAN/CAN FD 容忍节点 (\hyperref[fielddesc:TWAIFDFDE]{TWAIFD\_FDE} = 0) 时，若用户尝试发送 CAN FD 帧 (TX buffer 中的 FRAME\_FORMAT\_W[FDF\_BIT] = 1)，则模块将忽略该帧的类型，并发送 CAN 2.0 格式的帧。

当 \modulename{} 配置为经典 CAN/CAN FD 容忍节点时，\hyperref[fielddesc:TWAIFDNISOFD]{TWAIFD\_NISOFD} 寄存器无效。

根据 ISO11898-1:2015 的第 10.9.10 节，不应将 CAN FD 启用型实现配置为行为类似于 CAN FD 容忍型实现的模式。因此，用户应仅在评估或调试阶段使用该配置选项，并自行承担相关责任。

此外，根据原 CAN 2.0 规范，接收器应接受 R0 和 R1 位的任意取值；然而在 ISO11898-1:2015 中，规范要求经典 CAN 实现在检测到此类情况时发送错误帧。这一规范差异已在本实现中得到修正，以符合 ISO11898-1:2015 的最新要求。


\subsubsection{最小位时间/最大位速率}

系统时钟周期等于最小时间定额，从而影响 CAN 总线上可实现的最小位速率。\modulename{} 具有以下限制：

\begin{itemize}
    \item Phase\_Seg2 >= 2 个最小时间定额，适用于标称位速率和数据位速率。
    \item Sync\_Seg + Prop\_Seg + Phase\_Seg1 > 2 个最小时间定额，适用于标称位速率和数据位速率。
\end{itemize}

在满足上述条件时，每个位时间的长度最小可为 5 个时间定额。请注意，对于标称位速率而言，5 个时间定额虽然可行，但仍建议每个位时间至少使用 8 个时间定额，以获得更好的通信稳定性。
对于数据位速率，可以设置为 5 个时间定额。

例如，当标称位速率为 250 kbit/s、数据位速率为 1 Mbit/s 时，最小可用的系统时钟频率为 5 MHz。但请注意，这已是极限值，采样点位置的调整空间极小。因此，建议在该配置下使用至少 10 MHz 的系统时钟频率，确保有更安全的设计裕量。


\subsection{CAN 帧发送}
\label{sec:twai-can-frame-trans}

CAN 帧通过 \modulename{} 从 TX buffer 发起传输。\modulename{} 包含 4 个 TX buffer。当前 TX buffer 的数量可从 \hyperref[regdesc:TWAIFDTXCOMMANDTXTBINFOREG]{TWAIFD\_TX\_COMMAND\_TXTB\_INFO\_REG} 中读取。若存在 N 个缓冲区，则缓冲区编号固定为 1 至 N。每个 TX buffer 可以处于图 \ref{fig:twai-tx-buffer-state} 中描述的状态之一，从 \hyperref[regdesc:TWAIFDTXSTATUSREG]{TWAIFD\_TX\_STATUS\_REG} 中读取。软件能够通过 \hyperref[regdesc:TWAIFDTXCOMMANDTXTBINFOREG]{TWAIFD\_TX\_COMMAND\_TXTB\_INFO\_REG} 控制 TX buffer 状态，支持以下三种命令：

\begin{description}
    \item [set ready] 请求将 TX buffer 设置为 ready 状态。
    \item [set abort] 请求将 TX buffer 设置为 aborted 或 abort in progress 状态。
    \item [set empty] 请求将 TX buffer 设置为 empty 状态。
\end{description}

每个 TX buffer 可存储一帧 CAN 帧。整个 64 字节的 CAN FD 帧可写入单个 TX buffer。TX buffer 为只写寄存器（不支持读取回写入的帧），且仅在 TX buffer 处于 Empty、TX OK、TX failed 或 aborted 状态时可写入。CAN 帧写入 TX buffer 时使用的内存映射区域如下。

\begin{longtable}[]{@{}
  >{\raggedright\arraybackslash}p{(\linewidth - 2\tabcolsep) * \real{0.5000}}
  >{\raggedright\arraybackslash}p{(\linewidth - 2\tabcolsep) * \real{0.5000}}@{}}
\toprule\noalign{}
\begin{minipage}[b]{\linewidth}\raggedright
\bf{内存区域}
\end{minipage} & \begin{minipage}[b]{\linewidth}\raggedright
\bf{地址偏移}
\end{minipage} \\
\midrule\noalign{}
\endhead
\bottomrule\noalign{}
\endlastfoot
控制寄存器 & 0x000 \\
TX buffer 1 & 0x100 \\
TX buffer 2 & 0x200 \\
TX buffer 3 & 0x300 \\
TX buffer 4 & 0x400 \\
\end{longtable}

\begin{figure}[H]
\begin{centering}
\includegraphics{\modulefiles/figures/FSM_TXT_Buffer_user}
\par\end{centering}
\caption{TX buffer 状态}
\label{fig:twai-tx-buffer-state}
\end{figure}

当驱动程序将 CAN 帧写入 TX buffer 后，需向该 buffer 发出 set ready 命令以请求发送，此时 TX buffer 进入 ready 状态。在间歇字段的第三位采样到显性位（此时跳过 SOF），或总线空闲时，\modulename{} 将从 ready 进入 TX in progress 状态并开始发送。请注意，在定时发送的模式下行为有所不同，详见 \ref{sec:twai-func-time-trigger-trans}）。

当 \modulename{} 处于被动错误状态，且是上一帧的发送节点时，则会在连续发送前暂停 8 个位时间。当帧发送成功（即无仲裁丢失、无错误帧）时，TX buffer 转为 TX OK 状态。若出现错误帧或仲裁丢失，TX buffer 将重新进入 ready 状态，并在下一个间歇字段或总线空闲时重新尝试发送。

在总线监控模式 (\hyperref[fielddesc:TWAIFDBMM]{TWAIFD\_BMM} = 1) 或受限操作模式 (\hyperref[fielddesc:TWAIFDROM]{TWAIFD\_ROM} = 1) 下，即使发出 set ready 命令，\modulename{} 也将进入 TX failed 状态，不会尝试发送。


\subsubsection{TX buffer 选择}

如有多个 TX buffer 处于 ready 状态，\modulename{} 将从中选择优先级最高的 TX buffer 发送 CAN 帧，优先级由 \hyperref[regdesc:TWAIFDTXPRIORITYREG]{TWAIFD\_TX\_PRIORITY\_REG} 配置。如果两个 TX buffer 的优先级相同，则索引较小的 buffer 优先。发送流程如图\ref{fig:twai-tx-buffer-select} 所示。

\hyperref[regdesc:TWAIFDTXPRIORITYREG]{\hyperref[regdesc:TWAIFDTXPRIORITYREG]{TWAIFD\_TX\_PRIORITY\_REG}}{[}TX{*}P{]} 值越大表示 TX buffer {*} 优先级越高。例如，如果 \hyperref[regdesc:TWAIFDTXPRIORITYREG]{TWAIFD\_TX\_PRIORITY\_REG}{[}TX1P{]} = 2 且 \hyperref[regdesc:TWAIFDTXPRIORITYREG]{TWAIFD\_TX\_PRIORITY\_REG}{[}TX2P{]} = 5，则 TX buffer 2 的优先级为 5，优先级高于 TX buffer 1。两个 TX buffer 都处于 ready 状态时，\modulename{} 优先选择 TX buffer 2。

当 \hyperref[fielddesc:TWAIFDTXBBM]{TWAIFD\_TXBBM} = 1 时，候补 TX buffer 不通过其自身的 \hyperref[regdesc:TWAIFDTXPRIORITYREG]{TWAIFD\_TX\_PRIORITY\_REG}{[}TX{*}P{]} 配置优先级，而是由原 buffer 的对应字段配置。详见 \ref{sec:twai-func-tx-buffer-backup}。

\begin{figure}[H]
\begin{centering}
\includegraphics[scale=0.8]{\modulefiles/figures/TXT_Buffer_validation}
\par\end{centering}
\caption{TX buffer 选择}
\label{fig:twai-tx-buffer-select}
\end{figure}


\subsubsection{定时发送模式}
\label{sec:twai-func-time-trigger-trans}

\modulename{} 支持定时发送模式，对 \hyperref[fielddesc:TWAIFDTTTM]{TWAIFD\_TTTM} 写 1 可启用。在此模式下，仅当时间基准（参见章节 \ref{sec:twai-func-time-base}）的值达到 TX buffer 中 TIMESTAMP\_L\_W 和 TIMESTAMP\_U\_W 所存储的时间戳时，\modulename{} 才会尝试从优先级最高的 TX buffer 开始发送。时间基准是一个向上计数的无符号定时器，当其值达到设定的时间戳时，\modulename{} 将发送位于 ready 状态下优先级最高 TX buffer 中的帧，如图 \ref{fig:twai-time-trigger-trans} 所示。请注意，这并不意味着立即发送该帧，\modulename{} 仍需等待总线空闲。如果 TX buffer 处于 ready 状态，而时间基准尚未达到发送时间，\modulename{} 会等待直到条件满足。如果在此期间总线上已有其他节点开始发送，\modulename{} 会作为接收方参与通信。

\begin{figure}[H]
\begin{centering}
\includegraphics{\modulefiles/figures/TX_Time}
\par\end{centering}
\caption{定时发送}
\label{fig:twai-time-trigger-trans}
\end{figure}

若某帧应立即发送（即不依赖时间触发机制），则应将 TIMESTAMP\_L\_W、TIMESTAMP\_U\_W 配置为 0x00000000。请注意，定时发送仅考虑 ready 状态下优先级最高的 TX buffer。TX buffer 优先级始终高于发送时间，具体行为如下：

\begin{itemize}
    \item 如果 TX buffer A 的优先级高于 B，即使其发送时间晚于 B，\modulename{} 也会优先考虑 A。
    \item 如果优先级发生变化，导致 ready 状态下优先级最高的 buffer 改变，则\modulename{} 将选择发送新的最高优先级 buffer 中的帧。只要此前选中的 buffer 仍在等待时间基准达到发送时间，这一原则有效。一旦开始发送，将不再考虑优先级，不会切换 buffer 帧。
\end{itemize}

\subsubsection{发送 CAN 帧的类型}

\modulename{} 发送的 CAN 帧类型由 TX buffer 中 FRAME\_FORMAT\_W 字段和\modulename{} 的相关配置共同决定，如图 \ref{fig:twai-tx-frame-type} 所示。

\begin{figure}[H]
\begin{centering}
\includegraphics[scale=0.8]{\modulefiles/figures/TX_frame_flow}
\par\end{centering}
\caption{TX 帧类型}
\label{fig:twai-tx-frame-type}
\end{figure}

当 FRAME\_FORMAT\_W[FDF] = FD\_CAN 且\hyperref[fielddesc:TWAIFDFDE]{TWAIFD\_FDE} = 0 时，\modulename{} 发送 CAN 2.0 帧。如果此时 TX buffer 中为超过 8 字节的数据负载（CAN FD 格式），则超出 8 字节的部分不会发送。

当 FRAME\_FORMAT\_W[RTR] = RTR\_FRAME 且 FRAME\_FORMAT\_W[FDF] = FD\_CAN 时，将忽略 RTR 标志，\modulename{} 发送 CAN FD 帧（因为 CAN FD 协议不支持远程帧）。


\subsubsection{重发限制}

\modulename{} 可以限制单个 TX buffer 的重发次数，对 \hyperref[fielddesc:TWAIFDRTRLE]{TWAIFD\_RTRLE} 写 1 启。重发次数通过 \hyperref[fielddesc:TWAIFDRTRTH]{TWAIFD\_RTRTH} 配置，首次尝试发送 CAN 帧不计为重发。配置选项如表 \ref{tab:twai-retrans-counter-config} 所示。

\begin{table}[H]
    \centering
    \caption{重发限制配置}
    \label{tab:twai-retrans-counter-config}
    \begin{tabular}{ |l|l|m{9.2cm}| }
    \hline
    \rowcolor{lightgray}
    \textbf{\hyperref[fielddesc:TWAIFDRTRTH]{TWAIFD\_RTRTH}} & \textbf{\hyperref[fielddesc:TWAIFDRTRTH]{TWAIFD\_RTRTH}} & \textbf{行为}\\\hline
    - & 0 & 发送尝试没有任何限制，直到发送成功或进入离线。\\\hline
    0 & 1 & 仅尝试发送一次，失败后不重发，即一次性模式。\\\hline
    1～15 & 1 & 尝试发送帧 \hyperref[fielddesc:TWAIFDRTRTH]{TWAIFD\_RTRTH} + 1 次（首次发送 + \hyperref[fielddesc:TWAIFDRTRTH]{TWAIFD\_RTRTH} 次重发）。\\\hline
    \end{tabular}
\end{table}

如果单个 TX buffer 连续重发 \hyperref[fielddesc:TWAIFDRTRTH]{TWAIFD\_RTRTH} 次不成功（出现错误帧或仲裁丢失），则该 buffer 进入 TX failed 状态。如果两次发送之间 TX buffer 改变（例如优先级变化），内部重发计数器清零，新 TX buffer 的帧仍有 \hyperref[fielddesc:TWAIFDRTRTH]{TWAIFD\_RTRTH} + 1 次发送机会。如果再次使用原 TX buffer 发送，则重新开始计数，\modulename{} 将尝试发送 CAN 帧 \hyperref[fielddesc:TWAIFDRTRTH]{TWAIFD\_RTRTH} + 1 次。若当前的 TX buffer 进入 aborted 状态，内部重发计数器也将清零。如果 buffer 再次进入 ready 状态，\modulename{} 将尝试发送 \hyperref[fielddesc:TWAIFDRTRTH]{TWAIFD\_RTRTH} + 1 次。每个帧的发送尝试次数由一个内部计数器追踪，通过 \hyperref[fielddesc:TWAIFDRETRCTRVAL]{TWAIFD\_RETR\_CTR\_VAL} 读取。


\subsubsection{中止}

如果软件驱动程序先前通过 set ready 命令请求发送 CAN 帧，则可通过 set abort 命令中止该请求。若 TX buffer 接收 set abort 命令时仍处于 ready 状态（尚未开始发送），则会立即转为 aborted 状态。若 TX buffer 已处于 TX in progress 状态（已经开始发送），则转为 abort in progress 状态。在此情况下，一旦出现错误帧或仲裁丢失，将转换至 Aborted 状态。请注意，当 TX buffer 处于 abort in progress 状态时，如果当前发送成功，则进入 TX OK 状态；如果达到重发次数上限，则进入 TX failed 状态。

\subsubsection{TX buffer 离线行为}

当 \modulename{} 由于 TEC > 255 而离线时，TX buffer 可根据配置做出以下两种响应：

\begin{enumerate}
    \item 所有处于 ready、TX in progress 或 abort in progress 状态的 TX buffer 立即进入 TX failed 状态。对 \hyperref[fielddesc:TWAIFDTBFBO]{TWAIFD\_TBFBO} 写 1 启用，为默认配置。
    \item \modulename{} 离线时使用的 TX buffer 将按照常规错误帧处理，通过清零 \hyperref[fielddesc:TWAIFDTBFBO]{TWAIFD\_TBFBO} 启用。若未发出 set abort 命令且未达到重发上限，相关 buffer 会重新进入 ready 状态。在 \modulename{} 重新接入总线（详见章节 \ref{sec:twai-func-oper-mode}）后，将根据优先级规则继续尝试发送。此配置支持模块离线与重新接入的自动化，无需软件主动干预 TX buffer。
\end{enumerate}

配置流程详见章节 \ref{sec:twai-prog-can-frame-trans} \textit{\nameref{sec:twai-prog-can-frame-trans}}。


\subsection{CAN 帧接收}

\modulename{} 包含一个独立的 RX buffer，用于存储接收到的 CAN 帧。RX buffer 的大小是 32 位字的整数倍，通过 RX\_MEM\_INFO 读取。RX buffer 采用 FIFO 结构组织。当 CAN 总线上成功接收到一个 CAN 帧（未检测到错误帧）时，该 CAN 帧便存入 RX buffer。软件通过连续读取 \hyperref[regdesc:TWAIFDRXDATAREG]{TWAIFD\_RX\_DATA} 寄存器，从 RX buffer 中取出 CAN 帧。每次从 \hyperref[regdesc:TWAIFDRXDATAREG]{TWAIFD\_RX\_DATA} 寄存器读取，会从 RX buffer 读出一个字。RX buffer 支持以下两种操作模式：
\begin{itemize}
    \item 自动模式 - 读取 \hyperref[regdesc:TWAIFDRXDATAREG]{TWAIFD\_RX\_DATA} 时 RX FIFO 的读取指针自动递增，仅在通过 32 位访问方式读取 \hyperref[regdesc:TWAIFDRXDATAREG]{TWAIFD\_RX\_DATA} 时使用。此时，对 \hyperref[fielddesc:TWAIFDRXRPMV]{TWAIFD\_RXRPMV} 写 1 无效。
    \item 手动模式 - 读取 \hyperref[regdesc:TWAIFDRXDATAREG]{TWAIFD\_RX\_DATA} 时读取指针不会自动移动，需由软件对 \hyperref[fielddesc:TWAIFDRXRPMV]{TWAIFD\_RXRPMV} 写 1 手动增加。该模式适用于通过 8 位、16 位或 32 位访问方式读取 RX buffer，可借助多个小访问组合读取一个完整的字（如 4 次 8 位访问或 2 次 16 位访问）。
\end{itemize}

RX buffer 的模式由 \hyperref[fielddesc:TWAIFDRXBAM]{TWAIFD\_RXBAM} 配置。RX buffer 内的 CAN 帧格式可参考图 \ref{fig:twai-rx-buffer-frame}，该图展示了可能存储在 RX buffer 中的可选 CAN 帧格式（互斥）。RX buffer 中的每个 CAN 帧占用 4 到 20 个字，其大小计算公式为：

\begin{center}
RX 帧的字数 = 4 + ceil (数据字段长度/4)
\par\end{center}

\begin{figure}[H]
\begin{centering}
\includegraphics[scale=0.9]{\modulefiles/figures/RX_Buffer_layout}
\par\end{centering}
\caption{RX buffer}
\label{fig:twai-rx-buffer-frame}
\end{figure}


\subsubsection{帧数量}

RX buffer 内含一个帧计数器，统计 buffer 中存储的 CAN 帧数量，通过 \hyperref[fielddesc:TWAIFDRXFRC]{TWAIFD\_RXFRC} 读取。帧存入 RX buffer 时，计数器递增；从 RX buffer 读取到某个 CAN 帧的最后一个字时，计数器递减。


\subsubsection{RX buffer 内存}

RX buffer 内存提供以下状态信息：

\begin{itemize}
    \item 可用内存字数，通过 \hyperref[fielddesc:TWAIFDRXFREE]{TWAIFD\_RX\_FREE} 读取。
    \item 写指针位置，通过 \hyperref[fielddesc:TWAIFDRXWPP]{TWAIFD\_RX\_WPP} 读取。
    \item 读指针位置，通过 \hyperref[fielddesc:TWAIFDRXRPP]{TWAIFD\_RX\_RPP} 读取。
\end{itemize}

\subsubsection{RX buffer 状态}

没有存储任何 CAN 帧的 RX buffer 为空。在此状态下，\hyperref[fielddesc:TWAIFDRXE]{TWAIFD\_RXE} = 1。所有字都被 CAN 帧占用的 RX buffer 为满。在此状态下，\hyperref[fielddesc:TWAIFDRXF]{TWAIFD\_RXF} = 1。需要注意的是，即使 RX buffer 状态未满，例如还有 2 个空闲字，但可能此时最小的 CAN 帧也无法存入。

\subsubsection{溢出}

如果在接收 CAN 帧期间，RX buffer 没有足够的空闲空间存储当前接收的 CAN 帧，则会发生溢出，导致该帧被丢弃（RX buffer FIFO 溢出），同时溢出标志被置位。溢出标志为粘性标志（保持置位状态直至清除），可以通过 \hyperref[fielddesc:TWAIFDDOR]{TWAIFD\_DOR} 读取，通过对 \hyperref[fielddesc:TWAIFDCDO]{TWAIFD\_CDO} 写 1 清除。

\subsubsection{刷新}
\label{sec:twai-func-flush}

对 \hyperref[fielddesc:TWAIFDRRB]{TWAIFD\_RRB} 写 1 可以刷新 RX buffer。刷新将保留 buffer 中的内容（内存未被擦除），但会清零读写指针和帧计数器。此时，RX buffer 状态和未接收到任何帧相同。如果在 CAN 帧接收过程中执行此命令，当前正在接收的帧也会被丢弃。

\subsubsection{不一致保护}

从 RX buffer 读取 CAN 帧需要多次访问 \hyperref[regdesc:TWAIFDRXDATAREG]{TWAIFD\_RX\_DATA} 寄存器，每次读取会产生副作用，即 RX buffer 内的读取指针自动递增。如果在读取过程中发生错误（如总线错误或 ECC 错误），则已读取的数据可能丢失，且读取指针已经递增，导致软件驱动无法重新读取该字，从而软件与 RX buffer 状态不一致。软件驱动程序可以读取 \hyperref[fielddesc:TWAIFDRXMOF]{TWAIFD\_RXMOF} 来检测和恢复。当 RX buffer 的读取指针指向帧中间某个字（即下一次读取将返回除 FRAME\_FORMAT\_W 以外的字时），\hyperref[fielddesc:TWAIFDRXMOF]{TWAIFD\_RXMOF} 会被置 1。如果软件驱动出现不一致，可反复读取 \hyperref[regdesc:TWAIFDRXDATAREG]{TWAIFD\_RX\_DATA}，直至 \hyperref[fielddesc:TWAIFDRXMOF]{TWAIFD\_RXMOF} 返回 0。此时，\hyperref[regdesc:TWAIFDRXDATAREG]{TWAIFD\_RX\_DATA} 指针会定位到新帧的 FRAME\_FORMAT\_W 字，或者 RX buffer 为空（若在读取 RX buffer 中唯一帧时发生错误）。

\subsubsection{时间戳}

当 \modulename{} 接收到 CAN 帧时，会将其时间戳存储在 RX buffer 中的 TIMESTAMP\_L\_W 和 TIMESTAMP\_U\_W 字，时间戳为采样外部时间基准的值。接收帧的时间戳可以在帧起始位采样点，或帧结束后的第 6 位（依据 ISO11898-1:2015 标准中定义的 CAN 帧有效接收时间）进行采样。时间戳采样位置由 \hyperref[fielddesc:TWAIFDRTSOP]{TWAIFD\_RTSOP} 配置。

\subsubsection{帧过滤}

硬件过滤器用于筛选接收到的 CAN 帧。\modulename{} 包含两种类型的过滤器：屏蔽过滤器和范围过滤器，共有三个屏蔽过滤器实例 (A、B、C) 和一个范围过滤器实例。接收到的 CAN 帧通过任一过滤器（即过滤器结果按逻辑“或”计算），才会存入 RX buffer。帧过滤器在启用接受过滤器模式 (\hyperref[fielddesc:TWAIFDAFM]{TWAIFD\_AFM} = 1) 时生效，应用于接收到的帧标识符。\hyperref[fielddesc:TWAIFDAFM]{TWAIFD\_AFM} 仅可在 \modulename{} 禁用时修改 (\hyperref[fielddesc:TWAIFDENA]{TWAIFD\_ENA} = 0)。若禁用接受过滤器模式，过滤器功能无效，所有接收的 CAN 帧均会直接存入 RX buffer。

每个过滤器可选择配置为仅接受特定类型的 CAN 帧（CAN 2.0 或 CAN FD）及标识符类型（基本标识符或扩展标识符），通过 \hyperref[regdesc:TWAIFDFILTERCONTROLFILTERSTATUSREG]{TWAIFD\_FILTER\_CONTROL} 配置。如需禁用某过滤器，可将其对应的 \hyperref[regdesc:TWAIFDFILTERCONTROLFILTERSTATUSREG]{TWAIFD\_FILTER\_CONTROL} 寄存器全部字段清零。接受过滤器实例的内部工作机制如图 \ref{fig:twai-frame-filter-operation} 所示。


\begin{figure}[H]
\begin{centering}
\includegraphics[scale=0.7]{\modulefiles/figures/Frame_filters_flow}
\par\end{centering}
\caption{帧过滤器工作流程（{*} 以 A/B/C/R 类型过滤器为例）}
\label{fig:twai-frame-filter-operation}
\end{figure}


\subsubsection{屏蔽过滤器}

屏蔽过滤器会检查接收到的 CAN 帧标识符是否与 \hyperref[regdesc:TWAIFDFILTERAVALREG]{TWAIFD\_FILTER\_X\_VAL\_REG} 配置的值匹配（X 表示过滤器实例 A、B 或 C），仅比较 \hyperref[regdesc:TWAIFDFILTERAMASKREG]{TWAIFD\_FILTER\_X\_MASK\_REG} 指定的屏蔽值。

若仅用屏蔽过滤器过滤带基本标识符的帧，应配置 \hyperref[regdesc:TWAIFDFILTERAMASKREG]{TWAIFD\_FILTER\_X\_MASK\_REG{[}17:0{]}} 为 0b000000000000000000。


\subsubsection{范围过滤器}

范围过滤器判断接收帧的标识符是否落在下限 \hyperref[regdesc:TWAIFDFILTERRANLOWREG]{TWAIFD\_FILTER\_RAN\_LOW\_REG} 至上限 \hyperref[regdesc:TWAIFDFILTERRANHIGHREG]{TWAIFD\_FILTER\_RAN\_HIGH\_REG} 定义的十进制范围内。

若仅用屏蔽过滤器过滤带基本标识符的帧，应配置 \hyperref[regdesc:TWAIFDFILTERRANLOWREG]{TWAIFD\_FILTER\_RAN\_LOW\_REG{[}17:0{]}} 为 0b000000000000000000 和 \hyperref[regdesc:TWAIFDFILTERRANHIGHREG]{TWAIFD\_FILTER\_RAN\_HIGH\_REG{[}17:0{]}} 为 0b111111111111111111。

配置流程详见章节 \ref{sec:twai-prog-can-frame-recep} \textit{\nameref{sec:twai-prog-can-frame-recep}}。


\subsection{故障限制}

\modulename{} 的故障限制状态可以通过 \hyperref[regdesc:TWAIFDEWLERPFAULTSTATEREG]{TWAIFD\_EWL\_ERP\_FAULT\_STATE\_REG} 读取，状态转换见图 \ref{fig:twai-fault-confine}。故障限制计数器包括 RX 错误计数器 (REC) 和 TX 错误计数器 (TEC)，二者符合 ISO11898-1 标准定义，计数器值可分别从 REC 和 TEC 读取。此外，\modulename{} 提供区分标称位速率与数据位速率的错误计数器。标称位速率错误计数器 (ERR\_NORM) 通过 \hyperref[fielddesc:TWAIFDERRNORMVAL]{TWAIFD\_ERR\_NORM\_VAL} 读取，每检测到一次错误计数加 1。数据位速率错误计数器 (ERR\_FD) 通过 \hyperref[fielddesc:TWAIFDERRFDVAL]{TWAIFD\_ERR\_FD\_VAL} 读取，每检测到一次错误计数加 1。

所有四个计数器（REC、TEC、ERR\_NORM、ERR\_FD）均可由软件控制。由于此功能直接影响 \modulename{} 对 ISO11898-1 的合规性，因此仅在测试模式启用时 (\hyperref[fielddesc:TWAIFDTSTM]{TWAIFD\_TSTM} = 1) 时才可配置，所有计数器的值由 \hyperref[regdesc:TWAIFDCTRPRESREG]{TWAIFD\_CTR\_PRES\_REG} 配置。错误警告限制 (EWL) 和被动错误限制 (ERP) 由 \hyperref[fielddesc:TWAIFDEWL]{TWAIFD\_EWL} 和 \hyperref[fielddesc:TWAIFDERPLIMIT]{TWAIFD\_ERP\_LIMIT} 配置，默认值符合 ISO11898-1 标准。在测试模式下，EWL 和 ERP 字段支持写操作。

\begin{figure}[H]
\begin{centering}
\includegraphics{\modulefiles/figures/FSM_fault_confinement}
\par\end{centering}
\caption{故障限制}
\label{fig:twai-fault-confine}
\end{figure}

\subsection{容错}

\modulename{} 实现了以下容错机制：

\begin{itemize}
    \item RX buffer RAM 的奇偶校验保护
    \item TX buffer RAM 的奇偶校验保护
    \item TX buffer 备份模式 (\hyperref[fielddesc:TWAIFDTXBBM]{TWAIFD\_TXBBM})
\end{itemize}

运行上述机制必须满足以下条件：

\begin{itemize}
    \item 软件驱动程序需对 \hyperref[fielddesc:TWAIFDPCHKE]{TWAIFD\_PCHKE} 写 1 以启用奇偶校验错误检测。仅在 \hyperref[fielddesc:TWAIFDENA]{TWAIFD\_ENA} = 0 时可配置 \hyperref[fielddesc:TWAIFDPCHKE]{TWAIFD\_PCHKE}。
\end{itemize}

\subsubsection{RX buffer RAM 的奇偶校验保护}

\modulename{} 接收 CAN 帧并将其存储到 RX buffer RAM 时，会为 RX buffer RAM 的每个字添加一个奇偶校验位。软件驱动程序从 RX buffer RAM 读取帧时，可以通过读取 \hyperref[fielddesc:TWAIFDRXPE]{TWAIFD\_RXPE} 位检测奇偶校验错误。若 RX buffer RAM 中某字的奇偶校验位与计算结果不符，\modulename{} 会在每次读取 \hyperref[regdesc:TWAIFDRXDATAREG]{TWAIFD\_RX\_DATA} 寄存器时置位 \hyperref[fielddesc:TWAIFDRXPE]{TWAIFD\_RXPE}。

RX buffer RAM 中的伪单事件翻转 (SEU) 可能会修改 RX buffer RAM 中接收帧的 FRAME\_FORMAT\_W[DLC] 字，影响软件驱动程序看到的 RX 帧长度，进而引发 RX buffer 状态不一致。此时软件只读取了部分接收帧，后续从 RX buffer 读取的所有帧均被破坏。


\subsubsection{TX buffer RAM 的奇偶校验保护}
\label{sec:twai-parity-protect-tx-buffer}

软件将 CAN 帧存储到 TX buffer 时，\modulename{} 会为 TX buffer RAM 的每个字附加一个奇偶校验位。当 \modulename{} 尝试从 TX buffer 发送帧并检测到 TX buffer RAM 中的奇偶校验错误时，其行为如下：

\begin{enumerate}
    \item 如果 \modulename{} 在 FRAME\_FORMAT\_W、IDENTIFIER\_W、TIMESTAMP\_U\_W 或 TIMESTAMP\_L\_W 中检测到奇偶校验错误，则不会尝试发送该 CAN 帧。
    \item 如果上述 TX buffer 字无奇偶校验错误，则尝试发送该 CAN 帧。
    \item 如果在发送 CAN 帧时 \modulename{} 在 DATA\_1\_4\_W - DATA\_61\_64\_W 的任意字检测到奇偶校验错误，则会发送错误帧。
\end{enumerate}

当 \modulename{} 在 TX buffer RAM 中检测到步骤 1 - 3 所述的奇偶校验错误时，该 TX buffer 进入 parity error 状态，如图 \ref{fig:twai-tx-buffer-state} 所示，并置位 \hyperref[fielddesc:TWAIFDTXPE]{TWAIFD\_TXPE}。

若检测到奇偶校验错误，软件应在再次尝试发送前，将整个 CAN 帧重新写入 TX buffer。

对 \hyperref[fielddesc:TWAIFDCTXPE]{TWAIFD\_CTXPE} 写 1 可清除 \hyperref[fielddesc:TWAIFDTXPE]{TWAIFD\_TXPE} 位。

当 \hyperref[fielddesc:TWAIFDPCHKE]{TWAIFD\_PCHKE} = 0 时，\modulename{} 将忽略在 TX buffer 中的奇偶校验错误。此时不会配置 \hyperref[fielddesc:TWAIFDTXPE]{TWAIFD\_TXPE}，\hyperref[fielddesc:TWAIFDCTXPE]{TWAIFD\_CTXPE} 无效，TX buffer 不会进入 parity error 状态。



\modulename{} 不会检测 FRAME\_TEST\_W 中的奇偶校验错误。FRAME\_TEST\_W 用于故意破坏发送帧（例如测试 CAN 总线上的错误场景），该功能在需要奇偶校验保护的高可靠性容错应用中通常不适用。



\subsubsection{TX buffer 备份模式}\label{sec:twai-func-tx-buffer-backup}

\hyperref[fielddesc:TWAIFDTXBBM]{TWAIFD\_TXBBM} = 1 时，\modulename{} 进入 TX buffer 备份模式。在该模式下，具有相邻索引的两个 TX buffer 组成一对，如图 \ref{fig:twai-tx-buffer-pair} 所示。例如，\hyperref[fielddesc:TWAIFDTXTBUFFERCOUNT]{TWAIFD\_TXT\_BUFFER\_COUNT} = 8 时，\modulename{} 包含 8 个 TX buffer，共有 4 对 TX buffer，分别为 1-2、3-4、5-6、7-8。

\begin{figure}[H]
\begin{centering}
\includegraphics[scale=0.9]{\modulefiles/figures/TXT_buffer_pairs}
\par\end{centering}
\caption{TX buffer 对}\label{fig:twai-tx-buffer-pair}
\end{figure}

备份模式通过备份 TX buffer 为原 TX buffer 提供容错支持，buffer 对中索引较高的 buffer 作为备份，用于在原 buffer 发生奇偶校验错误时替代发送。工作流程如图 \ref{fig:twai-tx-buffer-backup-mode} 所示。

当 \hyperref[fielddesc:TWAIFDTXBBM]{TWAIFD\_TXBBM} = 1 且原 TX buffer 存在奇偶校验错误时，该 TX buffer 会进入 parity error 状态，\modulename{} 将尝试从对应的备份 TX buffer 发送帧（如 TX buffer 3 出错时，尝试从 TX buffer 4 发送）。若原 TX buffer 发送成功，其备份 TX buffer 则进入 aborted 状态，\modulename{} 不会发送备份帧。

当 \modulename{} 由于原 TX buffer 中的奇偶校验错误而从备份 TX buffer 发送帧时，如果在备份 TX buffer RAM 中也检测到奇偶校验错误，则会置位 \hyperref[fielddesc:TWAIFDTXDPE]{TWAIFD\_TXDPE}，表示存在双重奇偶校验错误。

在备份模式下，软件对 TX buffer 的控制与非备份模式存在差异：
\begin{itemize}
    \item TX buffer 对中的两个 TX buffer 的优先级相同，由原始 TX buffer 的 \hyperref[regdesc:TWAIFDTXPRIORITYREG]{TWAIFD\_TX\_PRIORITY\_REG}{[}TX{*}P{]} 决定。例如，TX buffer 1 和 2 的优先级由 \hyperref[regdesc:TWAIFDTXPRIORITYREG]{TWAIFD\_TX\_PRIORITY\_REG}{[}TX1P{]} 决定，\hyperref[regdesc:TWAIFDTXPRIORITYREG]{TWAIFD\_TX\_PRIORITY\_REG}{[}TX2P{]} 无效。
    \item 软件发出的命令将自动应用于每个原 TX buffer 及其对应的备份 TX buffer。例如，如果软件对 TX buffer 1 发出命令（TX\_COMMAND{[}TXB1{]} = 1），\modulename{} 会自动将其应用于 TX buffer 2。
\end{itemize}
软件存储相同的帧至 TX buffer 对中的两个 buffer，备份模式可在原 buffer 出现奇偶校验错误时，从备份 buffer 发送相应的帧。

\begin{figure}[H]
\begin{centering}
\includegraphics[scale=0.7]{\modulefiles/figures/TXT_Buffer_Backup_mode}
\par\end{centering}
\caption{TX buffer 备份模式工作流程}\label{fig:twai-tx-buffer-backup-mode}
\end{figure}

通过单独的内存访问将相同的帧存储到两个 TX buffer 是设计使然。\modulename{} 不会自动将帧复制到两个 TX buffer，以避免在写入过程中发生 SEU。如果发生 SEU，可能导致两个 TX buffer 中存储的帧已包含奇偶校验错误。

不必将完全相同的帧存储到 TX buffer 对中的两个 buffer。软件也可以仅将任意一帧存入备份 TX buffer，以便在原 TX buffer 出现奇偶校验错误时使用该备份帧进行发送。

软件应同时将 \hyperref[fielddesc:TWAIFDTXBBM]{TWAIFD\_TXBBM} 和 \hyperref[fielddesc:TWAIFDPCHKE]{TWAIFD\_PCHKE} 置 1。如果 \hyperref[fielddesc:TWAIFDTXBBM]{TWAIFD\_TXBBM} = 1，\hyperref[fielddesc:TWAIFDPCHKE]{TWAIFD\_PCHKE} = 0，\modulename{} 会忽略原 TX buffer 中的奇偶校验错误，且不会从备份 TX buffer 发送帧。

在 CAN 帧发送期间，如果 \modulename{} 在原 TX buffer 中检测到奇偶校验错误，且存在另一个优先级高于当前所用 TX buffer 对的 TX buffer 已因软件发出 set ready 命令而进入 ready 状态），\modulename{} 将在下一次发送时优先选择优先级更高的 TX buffer，忽略备份 TX buffer。

\modulename{} 仅在包含偶数个 TX buffer 时支持 TX buffer 备份模式。如果 TX buffer 数量为奇数，则存在一个没有备份 buffer 的 TX buffer。在此情况下，当 \hyperref[fielddesc:TWAIFDTXBBM]{TWAIFD\_TXBBM} = 1 时，软件不得使用这一无备份的原 TX buffer，否则 \modulename{} 行为将不可预测。


\subsection{特殊模式}

\subsubsection{回环模式}

回环模式下，\modulename{} 从任意 TX buffer 发送 CAN 帧后，若帧发送成功且通过帧过滤器，则将存储至 RX buffer。需注意，尽管\modulename{} 可将自身发送的帧存入 RX buffer，但其仍作为发送端，并不会确认自身的帧。要成功发送帧，须由 CAN 总线上的其他节点完成确认。回环模式通过对 \hyperref[fielddesc:TWAIFDILBP]{TWAIFD\_ILBP} 写 1 启用，仅在 \modulename{} 禁用 (\hyperref[fielddesc:TWAIFDENA]{TWAIFD\_ENA} = 0) 时可以修改。

\subsubsection{自检模式}

自检模式下，即使确认槽未收到显性位，\modulename{} 仍判定发送帧为有效。该模式可与回环模式同时启用，用于验证 \modulename{} 作为单节点在总线上的运行情况。自检模式通过对 \hyperref[fielddesc:TWAIFDSTM]{TWAIFD\_STM} 写 1 启用，仅在 \modulename{} 禁用 (\hyperref[fielddesc:TWAIFDENA]{TWAIFD\_ENA} = 0) 时可以修改。

\subsubsection{禁止确认模式}

禁止确认模式启用后，\modulename{} 作为 CAN 帧接收端时，即使接收到的 CRC 校验无误，也不会向确认槽发送显性位。禁止确认模式通过向 \hyperref[fielddesc:TWAIFDACF]{TWAIFD\_ACF} 写 1 启用，仅在 \modulename{} 禁用 (\hyperref[fielddesc:TWAIFDENA]{TWAIFD\_ENA} = 0) 时可以修改。


\subsubsection{总线监控模式}

总线监控模式下，\modulename{} 只接收、不发送任何帧。将 CAN 帧写入 TX buffer 并发出 set ready 命令时，不会发送该帧，且 TX buffer 会立即进入 TX failed 状态。在此模式下，\modulename{} 不在总线上发出任何显性位；当需发出显性位（如 ACK 或错误帧）时，显性位仅作为内部信号由 \modulename{} 接收，CAN 总线上其他节点无法检测到该字段。总线监控模式通过向 \hyperref[fielddesc:TWAIFDBMM]{TWAIFD\_BMM} 写 1 启用，仅在 \modulename{} 禁用 (\hyperref[fielddesc:TWAIFDENA]{TWAIFD\_ENA} = 0) 时可以修改。

\subsubsection{受限操作模式}

受限操作模式下，\modulename{} 能够接收 CAN 总线上的帧，但不发送任何帧。将 CAN 帧写入 TX buffer 并发出 set ready 命令时，不会发送该帧，且 TX buffer 会立即进入 TX failed 状态。在此模式下，\modulename{} 会确认有效帧，但不发送错误帧或过载帧。当检测到错误或过载时，\modulename{} 进入总线接入状态，并等待 11 个连续的隐性位。REC 与 TEC 计数器均不会被修改，\modulename{} 始终处于主动错误状态。受限操作模式通过向 \hyperref[fielddesc:TWAIFDROM]{TWAIFD\_ROM} 写 1 启用，仅在 \modulename{} 禁用 (\hyperref[fielddesc:TWAIFDENA]{TWAIFD\_ENA} = 0) 时可以修改。


\subsection{其他功能}

\subsubsection{错误代码捕获}

\modulename{} 配备了错误代码捕获寄存器，用于记录导致发送错误帧最后一次错误的类型和比特位置。检测到错误信号时，其采样点会实时更新该寄存器。用户可通过 \hyperref[regdesc:TWAIFDERRCAPTRETRCTRALCTSINFOREG]{TWAIFD\_ERR\_CAPT} 读取错误捕获信息。需要注意，CAN FD 标准未将错误类型定义为互斥（例如：某个填充位损坏时可能同时满足位错误和填充错误条件）。在这种情况下，错误代码捕获机制仅记录优先级最高的错误。错误类型的优先级如下，数字越小优先级越高（格式错误最高）：

\begin{center}
\begin{tabular}{|c|c|c|c|c|c|}
\hline
\textbf{优先级} \cellcolor{lightgray} & 1 & 2 & 3 & 4 & 5\tabularnewline
\hline
\textbf{错误类型} \cellcolor{lightgray} & 格式错误 & 位错误 & CRC 错误 & ACK 错误 & 填充错误\tabularnewline
\hline
\end{tabular}
\par\end{center}

填充 CAN FD 帧固定位时，发生的填充错误会在错误代码捕获寄存器中报告为格式错误。

上述错误优先级顺序有一个例外。如果在仲裁字段期间发送显性填充位时采样到隐性值，则此错误被捕获为填充错误，而不是位错误。


\subsubsection{仲裁丢失捕获}

\modulename{} 包含仲裁丢失捕获寄存器 (ALC)，用于记录 \modulename{} 最近一次在 CAN 仲裁字段丢失仲裁的比特位置。

\subsubsection{流量计数器}

\modulename{} 支持统计在 CAN 总线上发送和接收的 CAN 帧。每成功发送一个 CAN 帧，TX\_COUNTER 加 1。每成功接收一个 CAN 帧，RX\_COUNTER 加 1。通过向 \hyperref[fielddesc:TWAIFDTXFCRST]{TWAIFD\_TXFCRST}写 1 可清除 TX\_COUNTER，通过向 \hyperref[fielddesc:TWAIFDRXFCRST]{TWAIFD\_RXFCRST} 写 1 可清除 RX\_COUNTER。当\modulename{} 处于回环模式并且自身发送的帧存入 RX buffer 时，RX\_COUNTER 也会增加。流量计数器为 \modulename{} 的可选功能，可通过读取 \hyperref[fielddesc:TWAIFDSTCNT]{TWAIFD\_STCNT} 来判断其是否使能。

\subsubsection{调试寄存器}

\modulename{} 包含调试寄存器 \hyperref[regdesc:TWAIFDDEBUGREG]{TWAIFD\_DEBUG\_REG}，直接反映当前正在发送或接收的 CAN 帧的相关字段内容。


%%%%%%%%%%%%%%%%%%
%%% Interrupts %%%
%%%%%%%%%%%%%%%%%%

\section{中断}\label{sec:twai-interrupt}

\chipname{} 的 \modulename{} 可以生成以下中断信号并发送给 \hyperref[mod:intmtrx]{中断矩阵}。

    \begin{itemize}
        \item IRQ
        \item IRQ\_TIMER
    \end{itemize}

以下中断源可以中断源可以生成上述中断信号，具体取决于中断源：

\begin{itemize}
    \item \label{int:TWAIFDTXBHCIINT}TWAIFD\_TXBHCI\_INT：当 TX buffer 收到来自 CAN 核心的硬件命令，将 TX buffer 状态更改为 TX OK、error 或 aborted 时触发。
    \item \label{int:TWAIFDRBNEIINT}TWAIFD\_RBNEI\_INT：当 RX buffer 非空时触发。清除此中断且不通过 \hyperref[regdesc:TWAIFDRXDATAREG]{TWAIFD\_RX\_DATA} 读取 RX buffer 将重新激活中断。
    \item \label{int:TWAIFDBSIINT}TWAIFD\_BSI\_INT：当位速率发生变化时触发。
    \item \label{int:TWAIFDRXFIINT}TWAIFD\_RXFI\_INT：当 RX buffer 为满时触发。
    \item \label{int:TWAIFDOFIINT}TWAIFD\_OFI\_INT：当开始发送过载帧时触发。
    \item \label{int:TWAIFDBEIINT}TWAIFD\_BEI\_INT：当开始发送错误帧时触发。
    \item \label{int:TWAIFDALIINT}TWAIFD\_ALI\_INT：当仲裁丢失时触发。
     \item \label{int:TWAIFDFCSIINT}TWAIFD\_FCSI\_INT：当故障限制状态改变时触发，即节点状态从主动错误变为被动错误，从被动错误变为离线，从离线后重新接入或被动错误变为主动错误。
    \item \label{int:TWAIFDDOIINT}TWAIFD\_DOI\_INT：当数据溢出时触发。在清除此中断之前，必须清除 \hyperref[fielddesc:TWAIFDDOIINTST]{TWAIFD\_DOI\_INT\_ST} 以避免再次触发此中断。
    \item \label{int:TWAIFDEWLIINT}TWAIFD\_EWLI\_INT：当 TX 错误计数器 (TEC) 或 RX 错误计数器 (REC) 达到错误警告限制 (EWL) 时触发。当中断被清除且 REC 或 TEC 仍等于或高于 EWL 时，不会再次产生中断。
    \item \label{int:TWAIFDTXIINT}TWAIFD\_TXI\_INT：当发送帧时触发。
    \item \label{int:TWAIFDRXIINT}TWAIFD\_RXI\_INT：当接收帧时触发。
    \item \label{int:TWAIFDTIMEROVERFLOWINT}TWAIFD\_TIMER\_OVERFLOW\_INT：当定时器溢出时触发。
\end{itemize}

每个中断源都有一组通用的配置寄存器，详见章节 \hyperref[interrupt-config-registers]{\textit{中断配置寄存器}}。具体寄存器名称请查看 \ref{sec:twai-reg-summ} \textit{\nameref{sec:twai-reg-summ}}。


%%%%%%%%%%%%%%%%%%%%%%%%%%%%%%
%%% Programming Procedures %%%
%%%%%%%%%%%%%%%%%%%%%%%%%%%%%%

\section{配置示例}
\label{sec:twai-prog-procedures}


\subsection{500 kbit/4 Mbit 示例}
\label{sec:twai-prog-500-4-example}

在汽车行业中，CAN 总线常见的配置为：标称位速率 500 kbit/s，数据位速率 4 Mbit/s。对于 GPIO 选择无特定要求，收发器需选用支持 CAN FD 的型号（如 TD301MCANFD）。以下伪代码展示了在系统时钟源为 PLL\_F80M\_CLK 且采样点设为每个位的 80\% 时的典型配置方法：

\begin{lstlisting}[mathescape=true]
uint32 btr;
btr = (4 << 19);     // BRP = 4 (Tq = 4/80MHz = 50ns)
btr |= 14;           // PROP = 14
btr |= (17 << 7);    // SEG1 = 17
btr |= (8 << 13);   // SEG2 = 8
btr |= (2 << 27);    // SJW = 2
REG32_WR(TWAIFD_BTR_REG, btr); // (1+14+17+8)*50ns = 40*50ns = 2$\mu$s = 500kbit/s
uint32 btr_fd;
btr_fd = (1 << 19);     // BRP = 1 (Tq = 1/80MHz = 12.5ns)
btr_fd |= 3;           // PROP = 3
btr_fd |= (11 << 7);    // SEG1 = 11
btr_fd |= (5 << 13);   // SEG2 = 5
btr_fd |= (2 << 27);    // SJW = 2
REG32_WR(TWAIFD_BTR_FD_REG, btr_fd); // (1+3+11+5)*12.5ns = 20*12.5ns = 250ns = 4Mbit/s
\end{lstlisting}

按上述伪代码配置寄存器 \hyperref[regdesc:TWAIFDBTRREG]{TWAIFD\_BTR\_REG} 和 \hyperref[regdesc:TWAIFDBTRFDREG]{TWAIFD\_BTR\_FD\_REG}。如收发器选用 TD301MCANFD，建议配置 \hyperref[fielddesc:TWAIFDSSPOFFSET]{TWAIFD\_SSP\_OFFSET} 为 14。

\subsubsection{CAN 帧发送}\label{sec:twai-prog-can-frame-trans}

\paragraph{示例代码}
\begin{lstlisting}
#define CAN_FD_BASE TWAIFD_DEVICE_ID_VERSION_REG
#define TX_COMMAND_ADDR (CAN_FD_BASE + 0x74)
#define TXT_BUFFER_1_BASE (CAN_FD_BASE + 0x100)
/* 将 CAN 帧插入 TX buffer  1 */
uint32_t frame_format_word = 0;
frame_format_word |= 4;                               // DLC = 4
frame_format_word |= (1 << 7);                        // CAN FD 帧
frame_format_word |= (1 << 9);                        // 切换位速率
REG32_WR(TXT_BUFFER_1_BASE, frame_format_word);       // 存储帧格式字
uint32_t id_word = (55 << 18);                        // 标识符：55
REG32_WR(TXT_BUFFER_1_BASE + 0x4, id_word);           // 存储标识符字
REG32_WR(TXT_BUFFER_1_BASE + 0x8, 1000);
REG32_WR(TXT_BUFFER_1_BASE + 0xC, 0);                 // 在时间 1000 发送
REG32_WR(TXT_BUFFER_1_BASE + 0x10, 0xAABBCCDD);       // 数据：0xAA 0xBB 0xCC 0xDD
/* 发出 set ready 命令 */
uint32_t command = 0;
command |= 0x2;                                       // set ready 命令
command |= (1 << 8);                                  // 选择 TX buffer 1
REG32_WR(TX_COMMAND_ADDR, command);                   // 发出命令
\end{lstlisting}

向 \hyperref[fielddesc:TWAIFDENA]{TWAIFD\_ENA} 写 1 使能 \modulename{} 时，在接入 CAN 总线期间模块仍保持离线状态。如果此时向 TX buffer 下发 set ready 命令，并且 \hyperref[fielddesc:TWAIFDTBFBO]{TWAIFD\_TBFBO} = 1，TX buffer 会立即进入 aborted 状态。在向任意 TX buffer 发送 set ready 命令之前，软件应等待节点进入错误主动状态，可通过轮询 \hyperref[regdesc:TWAIFDEWLERPFAULTSTATEREG]{TWAIFD\_EWL\_ERP\_FAULT\_STATE\_REG} 或 FCS 中断判断。

TX buffer 未初始化或重置。因此，在发出 set ready 命令之前，软件必须先将有效的 CAN 帧填充至相应的 TX buffer。

\modulename{} 仅在检测到总线过载时主动发送过载帧，不会因内部原因自行发送过载帧。


\subsubsection{CAN 帧接收}\label{sec:twai-prog-can-frame-recep}

\paragraph{示例代码 - 自动模式下的帧接收（32 位访问方式）}
\begin{lstlisting}
#define CAN_FD_BASE TWAIFD_DEVICE_ID_VERSION_REG
#define RX_DATA_ADDR (CAN_FD_BASE + 0x6C)
#define RX_STATUS_ADDR (CAN_FD_BASE + 0x68)
/* 轮询 RX buffer 直到其中有帧 */
uint32_t rx_status;
do {
    rx_status = REG32_RD(RX_STATUS_ADDR);
} while ((rx_status \& 0x1) == 0)
/* 从 RX buffer 读取帧 */
uint8_t data[64];
uint32_t tmp;
uint32_t ffw = REG32_RD(RX_DATA_ADDR);
uint32_t id = REG32_RD(RX_DATA_ADDR);
uint32_t ts_l = REG32_RD(RX_DATA_ADDR);
uint32_t ts_h = REG32_RD(RX_DATA_ADDR);
uint32_t rwcnt = (ffw >> 11) \& 0x1F;
for(int i = 0; i < rwcnt; i++){
    tmp = REG32_RD(RX_DATA_ADDR);
    data[i*4] = tmp \& 0xFF;
    data[i*4+1] = (tmp >> 8) \& 0xFF;
    data[i*4+2] = (tmp >> 16) \& 0xFF;
    data[i*4+3] = (tmp >> 24) \& 0xFF;
}
\end{lstlisting}




%%%%%%%%%%%%%%%%%%%%%%%%
%%% Register Summary %%%
%%%%%%%%%%%%%%%%%%%%%%%%

\newpage
\hypertarget{twai-reg-summ}{}
\section{寄存器列表}
\label{sec:twai-reg-summ}

本小节的所有地址均为相对于 \modulename{} 基地址的地址偏移量（相对地址），具体基地址请见章节 \ref{mod:sysmem} \textit{\nameref{mod:sysmem}} 中的表 \ref{tab:sysmem-base-address}。

请查看章节 \hyperref[glossary-access-types]{\textit{寄存器的访问类型}}，了解“访问”列缩写的含义。

\subfile{\modulefiles/38-TWAI-reg-summ__CN}


%%%%%%%%%%%%%%%%%
%%% Registers %%%
%%%%%%%%%%%%%%%%%

\newpage
\section{寄存器}


本小节的所有地址均为相对于 \modulename{} 基地址的地址偏移量（相对地址），具体基地址请见章节 \ref{mod:sysmem} \textit{\nameref{mod:sysmem}} 中的表 \ref{tab:sysmem-base-address}。

关于保留 (reserved) 域的处理，请查看章节 \hyperref[programming-reserved-field]{\textit{如何配置寄存器的保留域}}。

\subfile{\modulefiles/38-TWAI-reg__CN}


%%%%%%%%%%%%%%%%%%%%%%%%%%%%%%%%%%%%%%%%%%%%%%%%%%%%%
%%% To paste template contents into your module's %%%
%%% root file, finish copying at the LINE ABOVE.  %%%
%%%%%%%%%%%%%%%%%%%%%%%%%%%%%%%%%%%%%%%%%%%%%%%%%%%%%


\end{document}
